\paragraph{Задачи}
\subparagraph{3.1. Асимптотическое поведение полиномов}
Пусть 
\[ p(n) = \sum_{i=0}^{d} a_i n^i, \]
где $a_d > 0$, представляет собой полином степени $d$ от $n$, и пусть $k$ является константой. Используя определения асимптотических обозначений, докажите следующие свойства.

\begin{itemize}
    \item[а.] Если $k \ge d$, то $p(n) = O(n^k)$.
    \item[б.] Если $k \le d$, то $p(n) = \Omega(n^k)$.
    \item[в.] Если $k = d$, то $p(n) = \Theta(n^k)$.
    \item[г.] Если $k > d$, то $p(n) = o(n^k)$.
    \item[д.] Если $k < d$, то $p(n) = \omega(n^k)$.
\end{itemize}

\subparagraph{3.2. Относительный асимптотический рост}
Для каждой пары приведенных в таблице выражений $(A, B)$ укажите, каким отношением $A$ связано с $B$: $O, o, \Omega, \omega$ или $\Theta$. Предполагается, что $k \ge 1, \epsilon > 0$ и $c > 1$ — константы. Ваш ответ должен выражаться таблицей, в каждой ячейке которой указано значение "Да" или "Нет".

\begin{center}
\renewcommand{\arraystretch}{1.5}
\begin{tabular}{cc|c|c|c|c|c}
\toprule
& $A$ & $B$ & $O$ & $o$ & $\Omega$ & $\omega$ & $\Theta$ \\ \midrule
\textbf{а.} & $\lg^k n$ & $n^\epsilon$ & & & & & \\ \hline
\textbf{б.} & $n^k$ & $c^n$ & & & & & \\ \hline
\textbf{в.} & $\sqrt{n}$ & $n^{\sin n}$ & & & & & \\ \hline
\textbf{г.} & $2^n$ & $2^{n/2}$ & & & & & \\ \hline
\textbf{д.} & $n^{\lg c}$ & $c^{\lg n}$ & & & & & \\ \hline
\textbf{е.} & $\lg(n!)$ & $\lg(n^n)$ & & & & & \\ 
\bottomrule
\end{tabular}
\end{center}

\subparagraph{3.3. Упорядочение по скорости асимптотического роста}
\begin{itemize}
    \item[а.] Расположите приведенные ниже функции по скорости их асимптотического роста, т.е. постройте такую последовательность функций $g_1, g_2, \dots, g_{30}$, что $g_1 = \Omega(g_2), g_2 = \Omega(g_3), \dots, g_{29} = \Omega(g_{30})$. Разбейте свой список на классы эквивалентности так, чтобы функции $f(n)$ и $g(n)$ находились в одном и том же классе тогда и только тогда, когда $f(n) = \Theta(g(n))$.
\end{itemize}

\begin{center}
\begin{tabular}{cccccc}
$\lg(\lg^* n)$ & $2^{\lg^* n}$ & $(\sqrt{2})^{\lg n}$ & $n^2$ & $n!$ & $(\lg n)!$ \\
$(3/2)^n$ & $n^3$ & $\lg^2 n$ & $\lg(n!)$ & $2^{2^n}$ & $n^{1/\lg n}$ \\
$\ln \ln n$ & $\lg^* n$ & $n \cdot 2^n$ & $n^{\lg \lg n}$ & $\ln n$ & $1$ \\
$2^{\lg n}$ & $(\lg n)^{\lg n}$ & $e^n$ & $4^{\lg n}$ & $(n+1)!$ & $\sqrt{\lg n}$ \\
$\lg^*(\lg n)$ & $2^{\sqrt{2 \lg n}}$ & $n$ & $2^n$ & $n \lg n$ & $2^{2^{n+1}}$
\end{tabular}
\end{center}

\begin{itemize}
    \item[б.] Приведите пример неотрицательной функции $f(n)$, такой, что для всех функций $g_i(n)$ из части (а) $f(n)$ не принадлежит ни множеству $O(g_i(n))$, ни множеству $\Omega(g_i(n))$.
\end{itemize}

\subparagraph{3.4. Свойства асимптотических обозначений}
Пусть $f(n)$ и $g(n)$ — асимптотически положительные функции. Докажите или опровергните справедливость каждого из приведенных ниже утверждений.

\begin{itemize}
    \item[а.] Из $f(n) = O(g(n))$ вытекает $g(n) = O(f(n))$.
    \item[б.] $f(n) + g(n) = \Theta(\min(f(n), g(n)))$.
    \item[в.] Из $f(n) = O(g(n))$ вытекает $\lg(f(n)) = O(\lg(g(n)))$, где $\lg(g(n)) \ge 1$ и $f(n) \ge 1$ для всех достаточно больших $n$.
    \item[г.] Из $f(n) = O(g(n))$ вытекает $2^{f(n)} = O(2^{g(n)})$.
    \item[д.] $f(n) = O((f(n))^2)$.
    \item[е.] Из $f(n) = O(g(n))$ вытекает $g(n) = \Omega(f(n))$.
    \item[ж.] $f(n) = \Theta(f(n/2))$.
    \item[з.] $f(n) + o(f(n)) = \Theta(f(n))$.
\end{itemize}

\subparagraph{3.5. Вариации определений $O$ и $\Omega$}
Некоторые авторы определяют $\Omega$ несколько иначе, чем это делали мы; давайте для такого альтернативного определения будем использовать символ $\overset{\infty}{\Omega}$ (читается как "омега бесконечность"). Будем говорить, что $f(n) = \overset{\infty}{\Omega}(g(n))$, если существует положительная константа $c$, такая, что $f(n) \ge cg(n) \ge 0$ для бесконечного количества целых чисел $n$.

\begin{itemize}
    \item[а.] Покажите, что для любых двух асимптотически неотрицательных функций $f(n)$ и $g(n)$ выполняется одно из соотношений $f(n) = O(g(n))$ и $f(n) = \overset{\infty}{\Omega}(g(n))$ (или они оба), в то время как при использовании $\Omega$ вместо $\overset{\infty}{\Omega}$ это утверждение ложно.
    \item[б.] Опишите потенциальные преимущества и недостатки применения $\overset{\infty}{\Omega}$ вместо $\Omega$ для характеристики времени работы программ.
\end{itemize}

Некоторые авторы определяют несколько иначе и $O$; для такого альтернативного определения будем использовать символ $O'$. Будем говорить, что $f(n) = O'(g(n))$ тогда и только тогда, когда $|f(n)| = O(g(n))$.

\begin{itemize}
    \item[в.] Что произойдет в теореме 3.1 с каждым из направлений "тогда и только тогда, когда", если заменить $O$ на $O'$, но оставить без изменений $\Omega$?
\end{itemize}

Некоторые авторы определяют $\tilde{O}$ (читается как "о с тильдой") для обозначения $O$, в котором игнорируются логарифмические множители:
\[ \tilde{O}(g(n)) = \{f(n) : \exists c, k, n_0 > 0 \text{ такие, что } 0 \le f(n) \le cg(n) \lg^k(n) \text{ для всех } n \ge n_0\}. \]

\begin{itemize}
    \item[г.] Определите $\tilde{\Omega}$ и $\tilde{\Theta}$ аналогичным образом. Докажите соответствующий аналог теоремы 3.1.
\end{itemize}

\subparagraph{3.6. Итерированные функции}
Оператор итерации $^*$, используемый в функции $\lg^*$, можно применить к любой монотонно возрастающей функции $f(n)$, определенной на множестве действительных чисел. Для заданной константы $c \in \mathbb{R}$ итерированная функция $f^*_c$ определяется следующим образом:
\[ f^*_c(n) = \min \{i \ge 0 : f^{(i)}(n) \le c\}. \]
Для каждой из приведенных далее функций $f(n)$ и констант $c$ дайте максимально точную оценку функции $f^*_c(n)$.

\begin{center}
\renewcommand{\arraystretch}{1.2}
\begin{tabular}{lcc|c}
\toprule
& $f(n)$ & $c$ & $f^*_c(n)$ \\ \midrule
\textbf{а.} & $n-1$ & $0$ & \\ \hline
\textbf{б.} & $\lg n$ & $1$ & \\ \hline
\textbf{в.} & $n/2$ & $1$ & \\ \hline
\textbf{г.} & $n/2$ & $2$ & \\ \hline
\textbf{д.} & $\sqrt{n}$ & $2$ & \\ \hline
\textbf{е.} & $\sqrt{n}$ & $1$ & \\ \hline
\textbf{ж.} & $n^{1/3}$ & $2$ & \\ \hline
\textbf{з.} & $n/\lg n$ & $2$ & \\ 
\bottomrule
\end{tabular}
\end{center}